\documentclass[12pt]{article}
\usepackage{resource/shortex}
\usepackage{mathtools,amssymb}
\usepackage{enumitem}
\usepackage{datetime}
\usepackage{fancyhdr}
\usepackage{ifthen}
\usepackage{pifont}
\usepackage{mathrsfs}
\usepackage[normalem]{ulem}
%\RequirePackage[usenames,dvipsnames]{color}


%% showing solutions 
\newboolean{showsols}
\setboolean{showsols}{false}
\newcommand{\showsolutions}{\setboolean{showsols}{true}}
\newlength{\solspace}
\setlength{\solspace}{16em}
\newcommand{\nosolspace}{\setlength{\solspace}{0em}}

\newcommand{\solution}[2]{\ifthenelse{\boolean{showsols}}{{\par\textcolor{red}{Rubric criterion: #1}\newline\textcolor{blue}{#2}}}{\vspace{\solspace}}}
\newcommand{\npforprint}{\ifthenelse{\boolean{showsols}}{}{\newpage}}
\newcommand{\vsforprint}[1]{\ifthenelse{\boolean{showsols}}{}{\vspace{#1}}}

%% grading
\newcommand{\grades}[1]{~\newline\noindent\framebox[\textwidth][c]{E\quad \quad \quad \quad S \quad \quad \quad \quad R \quad \quad \quad \quad N}}

\newcommand{\questiongrades}{\noindent\emph{Each item is marked $\checkmark$, $\checkmark\!-$, or \ding{55}:}\begin{center}\emph{$\checkmark$ = capable  $~~$ $\checkmark\!-$ = mostly capable  $~~$ \ding{55} = not yet}\end{center} }

\newcommand{\triAssessments}{$\checkmark\quad\checkmark\!-\quad$\ding{55}\quad}
\newcommand{\biAssessments}{$\checkmark\quad$\ding{55}\quad}

%% line for writing name and BUID
\newcommand{\nameBUID}{\noindent\textbf{Name:}\underline{\phantom{XXXXXXXXXXXXXXXXXXX}}\hfill\textbf{BUID:}\underline{\phantom{XXXXXXXXXXXXX}}\newline}


%% formatting 
%\allowdisplaybreaks[3]
\renewcommand{\headrulewidth}{.4mm} % header line width
\renewcommand{\arraystretch}{1.25}


%% page style 
\pagestyle{fancy}
\fancyhf{}
\rhead{Spring 2026}
\lhead{CAS MA 214: Applied Statistics}
%\rfoot{\thepage}

%--------------------
\begin{document}

\begin{center}
\textbf{\Large Project II Rubric} 
\end{center}

\grades

\subsection*{Rubric}

\questiongrades

\noindent
\section*{General Rubric}

\subsubsection*{Simulation Study Design (Part I)}
\normalsize
\begin{tabular}{lp{.75\textwidth}}
\toprule
\textbf{Mark} & \textbf{Description} \\
\hline
\triAssessments &\textbf{(1)} The base model and sampling procedures are correctly implemented and described with enough detail for replication. \\
\triAssessments &\textbf{(2)} The two chosen studies (A--D) are clearly identified and justified; the design is appropriate for the intended violation type and potential sources of bias are acknowledged. \\
\triAssessments &\textbf{(3)} Each violation scenario is accurately implemented in R with errors correctly generated. \\
\bottomrule
\end{tabular}
\vspace{1em}

\subsubsection*{Simulation Study Results (Part I)}
\normalsize
\begin{tabular}{lp{.75\textwidth}}
\toprule
\textbf{Mark} & \textbf{Description} \\
\hline
\triAssessments &\textbf{(1)} Coverage probabilities are reported in a well-organized table; bootstrap p-values for $H_0\colon \beta_1 = 0$ are compared to classical p-values; statistical tests are appropriate and multiple comparisons are interpreted carefully. \\
\triAssessments &\textbf{(2)} Average CI widths, bias of $\hat{\beta}_1$, and histograms of the sampling distribution of $\hat{\beta}_1$ are accurately reported; tables and figures are consistent with the narrative. \\
\triAssessments &\textbf{(3)} Results are interpreted with technically correct statistical reasoning; statistical and practical significance are distinguished.  \\
\bottomrule
\end{tabular}
\vspace{1em}

\subsubsection*{Real Data Analysis (Part II)}
\normalsize
\begin{tabular}{lp{.75\textwidth}}
\toprule
\textbf{Mark} & \textbf{Description} \\
\hline
\triAssessments &\textbf{(1)} The dataset is clearly described and cited; the predictor is motivated through EDA; the research question is concise and the OLS design is appropriate. \\
\triAssessments &\textbf{(2)} Classical and bootstrap 95\% CIs and p-values for $H_0\colon \beta_1 = 0$ are correctly computed and interpreted; agreement or disagreement between the two approaches is explicitly discussed. \\
\triAssessments &\textbf{(3)} Residual diagnostics are correctly produced and interpreted; OLS assumptions are assessed; the connection to the most analogous Part I scenario is argued with a well-reasoned conclusion about trust in the classical CI and p-value. \\
\bottomrule
\end{tabular}
\vspace{1em}

\subsubsection*{Validity of Conclusions (Part III \& Conclusion)}
\normalsize
\begin{tabular}{lp{.75\textwidth}}
\toprule
\textbf{Mark} & \textbf{Description} \\
\hline
\triAssessments &\textbf{(1)} Conclusions follow logically from the results; interpretations are accurate and not overstated \\
\triAssessments &\textbf{(2)} Practical recommendations are grounded in simulation evidence \\
\bottomrule
\end{tabular}
\vspace{1em}

\subsubsection*{Report Quality, Replicability, and R Code}
\normalsize
\begin{tabular}{lp{.75\textwidth}}
\toprule
\textbf{Mark} & \textbf{Description} \\
\hline
\triAssessments &\textbf{(1)} Report is well-structured; all sources are cited and attributed; no issues of plagiarism or missing references. \\
\triAssessments &\textbf{(2)} Visuals are clear, labeled, and support the narrative; text is fully consistent with all tables and figures. \\
\triAssessments &\textbf{(3)} R script is clean, reproducible, and commented; \texttt{set.seed()} is used throughout; the analysis is replicable from the report alone. \\
\bottomrule
\end{tabular}
\vspace{1em}

\section*{Individual Rubric}
\normalsize
\begin{tabular}{lp{.75\textwidth}}
\toprule
\textbf{Mark} & \textbf{Description} \\
\hline
\triAssessments &\textbf{(1)} Completed assigned individual tasks thoroughly. \\
\triAssessments &\textbf{(2)} Actively participated in group collaboration and communication. \\
\bottomrule
\end{tabular}
\vspace{1em}

\newpage
\subsection*{Grading Standards}

\begin{itemize}
\item \textbf{E: Excellent}: Satisfactory with exceptional creativity or insight
\item \textbf{S: Satisfactory}: Any combination of $\checkmark$ and $\checkmark\!-$ 
\item \textbf{R: Revisions Needed}: Any \ding{55} on any criterion
\item \textbf{N: Not Assessable}: Cannot grade due to too many missing components
\end{itemize}

\end{document}